\section{Discussion}\label{sec:discussion}
\subsection{What the evidence supports}
\emph{(1) Fusion helps when both modalities carry independent evidence.} In the controlled benchmark every fusion model beat the better single modality by about \gainUni{} points (CAFN, $n{=}\nGainUni$ paired seeds, $p\gainUniP$), and the sweep shows that this advantage is largest at intermediate informativeness and disappears when one modality dominates. This is the behaviour the multimodal-fusion literature assumes \cite{calhoun2016,baltrusaitis2019}, now measured under controlled conditions.

\emph{(2) The fusion mechanism mattered little.} Gated bidirectional cross-attention was statistically indistinguishable from concatenation (\dCat{} points, $p\pCat$) and element-wise sum (\dSum{}, $p\pSum$), and none of the seven single-component ablations produced a reliable change (range \ablMin{} to \ablMax{} points; $\ablSig$ of $\ablN$ nominally significant). CAMF reports that cross-attention beats concatenation and summation by roughly 1--3 accuracy points (arXiv version) on real schizophrenia data \cite{zhou2024camf}, a gap of the order of our seed-to-seed standard deviation (\sdAll{} points); we can neither confirm nor refute it, but at this model and sample size we do not find that the extra machinery buys accuracy. What cross-attention does provide is a by-product that concatenation lacks: attention maps in both directions (and a P100 estimator) that can be inspected, which is the stated clinical motivation. The faithfulness of those maps is not established here.

\emph{(3) On the real cohorts the signal is weak, and fusion did not help.} Only the spectral stream carried a signal (AUC 0.56--0.62, with the neural spectrum-only model highest), the MRI stream was at chance, the fused models were at chance, and nothing transferred between cohorts. We regard this as the expected outcome rather than a failure of the pipeline: 95 subjects, two unrelated disorders pooled under one label, different sensors, and a model of about 2\,M parameters that memorises its $\sim$67 training subjects (Fig.~\ref{fig:realcurves}).

\subsection{Why the real-data result is null}
Several factors are confounded and cannot be separated with these data. (i)~\emph{Sample size}: 95 subjects give AUC intervals about $\pm0.12$ wide. (ii)~\emph{Label heterogeneity}: depression and zoster neuralgia are different conditions, so a single ``patient'' class is not one concept. (iii)~\emph{Surrogate signal}: the resting spectrum is not a VEP and no timing biomarker such as P100 latency is available. (iv)~\emph{Preprocessing}: skull-stripping is still a threshold mask. In a single-seed diagnostic on the synthetic generator, switching the MRI pipeline from a percentile mask with min--max scaling to an Otsu mask with foreground z-scoring dropped the MRI-only accuracy from 73\% to 41\%, so real-MRI preprocessing is a first-order variable. (v)~\emph{Cohort effects}: the spectrum identifies the cohort perfectly, so any pooled model can exploit site; the within-cohort permutation test and leave-one-cohort-out evaluation exist to expose that.

\subsection{Pitfalls we hit and fixed}
These are of independent practical interest: restoring the checkpoint with the lowest \emph{validation loss} restored an untrained epoch-1 model on short runs (the early validation loss spikes while batch-norm statistics settle), which is why selection is now on validation balanced accuracy; a centre crop of a whole-head scan to $64^3$ sees only a sliver of the head, which is why the real-MRI path crops to the foreground bounding box; PyTorch's adaptive average pooling has a non-deterministic CUDA backward and no Apple MPS kernel, which we replaced with an equivalent matrix-product implementation; and running on a second platform exposed three defects hidden on the first (Section~\ref{sec:results}).

\subsection{Limitations}
(i)~Study A is synthetic: the modalities are independent given the label by construction, which favours fusion; real data may behave differently. (ii)~The number of seeds is small (at most \nSeedsMax{} per configuration, two per sweep point), so only large effects are detectable. (iii)~The real cohorts are not eye-disease data, the label is pan-disorder, and the 1-D stream is a spectrum, not a VEP; no claim about glaucoma, optic neuritis or CVI follows. (iv)~Skull-stripping and harmonisation across scanners and sensors are crude. (v)~Neural models on real data were cross-validated once and have no permutation test. (vi)~Explainability outputs were not validated against ground truth. (vii)~Hyper-parameters of the real-data study were fixed in advance rather than tuned, so a better-tuned model might do somewhat better; we chose not to tune on 95 subjects.

\subsection{Contribution in context}
Relative to \cite{zhou2024camf,dhar2025drifa}, which report fusion gains on benchmark data, the contribution here is methodological rather than a new state of the art: a fusion architecture for \emph{heterogeneous} 3D-image + 1-D-signal pairs, a controlled benchmark that isolates the effect of fusion from the effect of architecture, a reproducible one-command evaluation protocol with cross-cohort and permutation-based statistics, and an explicit account of what does and does not hold.

\section{Conclusion and Future Work}\label{sec:conclusion}
We presented CAFN, a gated bidirectional cross-attention network for paired 3D MRI and 1D electrophysiology, and evaluated it in two ways. In a controlled synthetic benchmark fusion improved accuracy by about \gainUni{} points over the better single modality, while the choice among cross-attention, concatenation and summation, and every individual design component, made no reliable difference ($\accFull\%$ for CAFN versus $\accCat\%$ and $\accSum\%$). On two public real cohorts (95 subjects, pan-disorder label, T1w MRI plus resting MEG/EEG) neural fusion was at chance, a weak spectral signal (AUC about 0.56--0.62) was detectable, and nothing generalised across cohorts. The whole pipeline, from download to test statistics, runs from one command with a fixed seed. The next step is data, not architecture: a labelled, paired MRI--VEP cohort with patient-level splits; pre-trained or self-supervised encoders and strong regularisation for small samples; harmonisation across sites; validated skull-stripping; and checks of the attention and Grad-CAM maps against clinician-marked regions or the known lesion masks of the synthetic generator.
